%==============================================================================
% APPENDIX A --- DATASET DETAILS AND CLASS TAXONOMY
%==============================================================================

\chapter{Dataset Details and Class Taxonomy}
\label{app:datasets}

\section{Dataset Sources}
\label{sec:app_dataset_sources}

The plant disease classification dataset used in \nabta was constructed by merging nine public datasets. These sources were selected to increase plant-species coverage, include both controlled and field-acquired images, and improve the diversity of disease symptoms represented during model training. The source datasets are summarized in Table~\ref{tab:dataset_sources}.

{\scriptsize
\begin{longtable}{@{}L{2.8cm} L{1.8cm} L{5.0cm} L{1.7cm} L{3.7cm}@{}}
\caption{Public Dataset Sources Used for Plant Disease Classification}
\label{tab:dataset_sources}\\
\toprule
\textbf{Dataset Name} & \textbf{Source} & \textbf{Description} & \textbf{Reference} & \textbf{Access Link} \\
\midrule
\endfirsthead
\caption[]{Public Dataset Sources Used for Plant Disease Classification (continued)}\\
\toprule
\textbf{Dataset Name} & \textbf{Source} & \textbf{Description} & \textbf{Reference} & \textbf{Access Link} \\
\midrule
\endhead
\midrule
\multicolumn{5}{r}{\textit{Continued on next page}}\\
\endfoot
\bottomrule
\endlastfoot
PlantVillage & Kaggle & Benchmark collection of healthy and diseased crop leaf images commonly used for plant disease classification. & \cite{dataset_plantvillage_kaggle} & \href{https://www.kaggle.com/datasets/mohitsingh1804/plantvillage}{Kaggle link} \\
Plant Diseases Training Dataset & Kaggle & Multi-crop plant disease image dataset used to extend disease coverage and class diversity. & \cite{dataset_plant_diseases_training} & \href{https://www.kaggle.com/datasets/nirmalsankalana/plant-diseases-training-dataset}{Kaggle link} \\
Bangladesh Plant Disease Dataset & Mendeley Data & Field-oriented plant disease dataset containing crop images collected under practical agricultural conditions. & \cite{dataset_bangladesh_plant_disease} & \href{https://data.mendeley.com/datasets/n67gctmjyj/3}{Mendeley link} \\
Eggplant Leaf Disease Dataset & Kaggle & Eggplant leaf disease images representing multiple visual symptoms and healthy samples. & \cite{dataset_eggplant_leaf_kaggle} & \href{https://www.kaggle.com/datasets/researchforus/eggplant-leaf-disease-dataset}{Kaggle link} \\
Eggplant Disease Dataset & Mendeley Data & Additional eggplant disease images used to improve robustness for eggplant-related classes. & \cite{dataset_eggplant_mendeley} & \href{https://data.mendeley.com/datasets/pvsv534ccg/1}{Mendeley link} \\
Rice Leaf Diseases Dataset & Kaggle & Rice leaf disease dataset covering major rice disease symptoms and healthy conditions. & \cite{dataset_rice_leaf_kaggle} & \href{https://www.kaggle.com/datasets/vbookshelf/rice-leaf-diseases}{Kaggle link} \\
Cotton Disease Dataset & Kaggle & Cotton leaf images containing healthy and diseased samples collected for crop disease recognition. & \cite{dataset_cotton_kaggle} & \href{https://www.kaggle.com/datasets/janmejaybhoi/cotton-disease-dataset}{Kaggle link} \\
Mango Leaf Disease Dataset & Kaggle & Mango leaf disease dataset covering several mango disease categories and healthy leaves. & \cite{dataset_mango_kaggle} & \href{https://www.kaggle.com/datasets/aryashah2k/mango-leaf-disease-dataset}{Kaggle link} \\
Ricey Leaf Disease Dataset & Mendeley Data & Rice leaf disease dataset collected for real-world rice disease classification research. & \cite{dataset_ricey_mendeley} & \href{https://data.mendeley.com/datasets/t46kkgh2yw/1}{Mendeley link} \\
\end{longtable}
}

\section{Class Distribution}
\label{sec:app_class_distribution}

The final classification taxonomy contains 102 output classes grouped by plant species. Table~\ref{tab:class_distribution} summarizes the number of healthy and diseased classes per species. Overall, the taxonomy contains 24 healthy classes and 78 disease classes.

\begin{longtable}{@{}L{4.4cm} C{2.6cm} C{2.6cm} C{2.6cm}@{}}
\caption{Class Distribution by Plant Species}
\label{tab:class_distribution}\\
\toprule
\textbf{Plant Species} & \textbf{Healthy Classes} & \textbf{Disease Classes} & \textbf{Total Classes} \\
\midrule
\endfirsthead
\caption[]{Class Distribution by Plant Species (continued)}\\
\toprule
\textbf{Plant Species} & \textbf{Healthy Classes} & \textbf{Disease Classes} & \textbf{Total Classes} \\
\midrule
\endhead
\bottomrule
\endlastfoot
Apple & 1 & 6 & 7 \\
Bitter Gourd & 1 & 3 & 4 \\
Blueberry & 1 & 0 & 1 \\
Bottle Gourd & 1 & 2 & 3 \\
Cassava & 1 & 3 & 4 \\
Cauliflower & 1 & 2 & 3 \\
Cherry (Including Sour) & 1 & 1 & 2 \\
Coffee & 1 & 1 & 2 \\
Corn (Maize) & 1 & 3 & 4 \\
Cotton & 1 & 1 & 2 \\
Cucumber & 1 & 2 & 3 \\
Eggplant & 1 & 8 & 9 \\
Grape & 1 & 3 & 4 \\
Mango & 1 & 7 & 8 \\
Orange & 0 & 1 & 1 \\
Peach & 1 & 1 & 2 \\
Pepper (Bell) & 1 & 1 & 2 \\
Potato & 1 & 7 & 8 \\
Raspberry & 1 & 0 & 1 \\
Rice & 1 & 9 & 10 \\
Rose & 1 & 2 & 3 \\
Soybean & 1 & 0 & 1 \\
Squash & 0 & 1 & 1 \\
Strawberry & 1 & 1 & 2 \\
Tomato & 1 & 10 & 11 \\
Watermelon & 1 & 3 & 4 \\
\midrule
\textbf{Total} & \textbf{24} & \textbf{78} & \textbf{102} \\
\end{longtable}

The distribution in Table~\ref{tab:class_distribution} shows that the taxonomy intentionally emphasizes disease categories while preserving healthy reference classes for most supported plant species. Species such as tomato, rice, eggplant, mango, and potato contain the largest disease coverage because they are represented by multiple public datasets and contain several visually distinct disease symptoms.

\section{Complete Class Taxonomy}
\label{sec:app_classes}

Table~\ref{tab:class_taxonomy} lists the complete set of 102 model output classes. The class identifiers correspond to the output indices produced by the plant disease classification model.

\begin{longtable}{@{}C{1.6cm} L{4.2cm} L{8.1cm}@{}}
\caption{Complete 102-Class Plant Disease Classification Taxonomy}
\label{tab:class_taxonomy}\\
\toprule
\textbf{Class ID} & \textbf{Plant Species} & \textbf{Disease / Condition} \\
\midrule
\endfirsthead
\caption[]{Complete 102-Class Plant Disease Classification Taxonomy (continued)}\\
\toprule
\textbf{Class ID} & \textbf{Plant Species} & \textbf{Disease / Condition} \\
\midrule
\endhead
\midrule
\multicolumn{3}{r}{\textit{Continued on next page}}\\
\endfoot
\bottomrule
\endlastfoot
0 & Apple & Apple Scab \\
1 & Apple & Black Rot \\
2 & Apple & Cedar Apple Rust \\
3 & Apple & Alternaria Leaf Spot \\
4 & Apple & Brown Spot \\
5 & Apple & Gray Spot \\
6 & Apple & Healthy \\
7 & Bitter Gourd & Downy Mildew \\
8 & Bitter Gourd & Fusarium Wilt \\
9 & Bitter Gourd & Healthy \\
10 & Bitter Gourd & Mosaic Virus \\
11 & Blueberry & Healthy \\
12 & Bottle Gourd & Anthracnose \\
13 & Bottle Gourd & Downy Mildew \\
14 & Bottle Gourd & Healthy \\
15 & Cassava & Brown Streak Disease \\
16 & Cassava & Green Mottle \\
17 & Cassava & Healthy \\
18 & Cassava & Mosaic Disease \\
19 & Cauliflower & Black Rot \\
20 & Cauliflower & Downy Mildew \\
21 & Cauliflower & Healthy \\
22 & Cherry (Including Sour) & Powdery Mildew \\
23 & Cherry (Including Sour) & Healthy \\
24 & Coffee & Healthy \\
25 & Coffee & Rust \\
26 & Corn (Maize) & Cercospora Leaf Spot (Gray Leaf Spot) \\
27 & Corn (Maize) & Common Rust \\
28 & Corn (Maize) & Northern Leaf Blight \\
29 & Corn (Maize) & Healthy \\
30 & Cotton & Diseased Leaf \\
31 & Cotton & Healthy Leaf \\
32 & Cucumber & Anthracnose \\
33 & Cucumber & Downy Mildew \\
34 & Cucumber & Healthy \\
35 & Eggplant & Cercospora Leaf Spot \\
36 & Eggplant & Hadda Beetles \\
37 & Eggplant & Healthy \\
38 & Eggplant & Insect Pest Disease \\
39 & Eggplant & Leaf Spot \\
40 & Eggplant & Mosaic Virus \\
41 & Eggplant & Phomopsis Blight \\
42 & Eggplant & Tobacco Caterpillar \\
43 & Eggplant & Wilt \\
44 & Grape & Black Rot \\
45 & Grape & Esca (Black Measles) \\
46 & Grape & Leaf Blight (Isariopsis Leaf Spot) \\
47 & Grape & Healthy \\
48 & Mango & Anthracnose \\
49 & Mango & Bacterial Canker \\
50 & Mango & Cutting Weevil \\
51 & Mango & Die Back \\
52 & Mango & Gall Midge \\
53 & Mango & Healthy \\
54 & Mango & Powdery Mildew \\
55 & Mango & Sooty Mould \\
56 & Orange & Huanglongbing (Citrus Greening) \\
57 & Peach & Bacterial Spot \\
58 & Peach & Healthy \\
59 & Pepper (Bell) & Bacterial Spot \\
60 & Pepper (Bell) & Healthy \\
61 & Potato & Early Blight \\
62 & Potato & Late Blight \\
63 & Potato & Bacterial Wilt \\
64 & Potato & Healthy \\
65 & Potato & Leafroll Virus \\
66 & Potato & Mosaic Virus \\
67 & Potato & Pests \\
68 & Potato & Phytophthora \\
69 & Raspberry & Healthy \\
70 & Rice & Bacterial Blight \\
71 & Rice & Blast \\
72 & Rice & Brown Spot \\
73 & Rice & Healthy \\
74 & Rice & Leaf Scald \\
75 & Rice & Leaf Smut \\
76 & Rice & Narrow Brown Leaf Spot \\
77 & Rice & Rice Hispa \\
78 & Rice & Sheath Blight \\
79 & Rice & Tungro \\
80 & Rose & Healthy \\
81 & Rose & Rust \\
82 & Rose & Slug Sawfly \\
83 & Soybean & Healthy \\
84 & Squash & Powdery Mildew \\
85 & Strawberry & Leaf Scorch \\
86 & Strawberry & Healthy \\
87 & Tomato & Bacterial Spot \\
88 & Tomato & Early Blight \\
89 & Tomato & Late Blight \\
90 & Tomato & Leaf Mold \\
91 & Tomato & Septoria Leaf Spot \\
92 & Tomato & Spider Mites (Two-Spotted Spider Mite) \\
93 & Tomato & Target Spot \\
94 & Tomato & Tomato Yellow Leaf Curl Virus \\
95 & Tomato & Tomato Mosaic Virus \\
96 & Tomato & Healthy \\
97 & Tomato & Spotted Wilt \\
98 & Watermelon & Anthracnose \\
99 & Watermelon & Downy Mildew \\
100 & Watermelon & Healthy \\
101 & Watermelon & Mosaic Virus \\
\end{longtable}
